Our %\sout{mathematical} 
study of t-SNE has established in considerable generality that one cannot infer the \textit{degree} of cluster salience or the extremity of outliers from a t-SNE plot, see Theorems \ref{thm:unclustHammer}, \ref{thm:perturbhammer}, and \ref{thm:outlier_abs}. The proofs and intuitions behind these statements guided us to surprising empirical observations that one cannot even infer the \textit{existence} of clusters or outliers in some cases. In particular, the injection of a small subset of adversarially chosen points can largely mask the cluster structure, while sizable injections of outlier points are reliably masked within the cluster structure, 
see Figures \ref{fig:one_pt_perturb}, \ref{fig:outliers1}, \ref{fig:outliers_real_world}, and \ref{fig:bbc_appendix}. %Further work should seek to formalize these latter set of empirical observations. C


%These behaviors are more pronounced as data becomes more high-dimensional: the damage done by an adversarially chosen poison point is largest on high-aspect-ratio data, while more faraway outliers can fly under the radar in the high-dimensional regime.   %t-SNE on high-dimensional data is susceptible to certain adversarial attacks and 

%Broadly speaking, w
We have identified two properties of t-SNE that give rise to these idiosyncratic behaviors: (1) additive invariance with respect to the squared interpoint distances, and (2) the asymmetry between the input and output affinity matrices. % (in that the former probes nearest-neighbors and the latter probes radius-neighbors)
 While we have uncovered significant failure modes that arise from these properties, we cannot completely rule out their utility. {For instance, though additive invariance may lead to exaggerated clusters, prior work has indicated that it may be useful in filtering out high-dimensional noise that may be present in the input, see e.g.\ Figure \ref{fig:higher_dim_tighter_clusters}, and \citet{lee2011, lee2014, Karoui2010, Kaouri2015, Landa2023}}.
 %Additive invariance, while brittle under certain adversarial perturbations, may be robust to certain random perturbations. Indeed, adding random noise to a dataset is approximately equivalent to adding a constant to the interpoint distances due to concentration of measure. Additive invariance effectively allows t-SNE to ignore such noise, see Figure \ref{fig:higher_dim_tighter_clusters}. We believe this phenomenon is worthy of further study.


 
 %while additive invariance is susceptible to adversarial perturbations, it is perhaps very robust to random perurbation
 
% additive invariance has desirable denoising properties, considering the concentration effects of high-dimensional noise (see Figure \ref{fig:higher_dim_tighter_clusters}).  Indeed, Gaussian noise 

%The holy grail of d
%The gold standard result for data visualization is a guarantee that clustered output implies clustered input in a suitable sense.



%t-SNE belongs to a larger collection of practical data visualization techniques \citep{mcinnes2018umap, jacomy2014forceatlas2, tang2016visualizing, amid2019trimap}. Beyond characterizing the intrinsic strengths and limitations of t-SNE, it would be instructive to taxonomize this zoo of techniques.

t-SNE belongs to a wide selection of data visualization techniques that are yet to be understood fully \citep{mcinnes2018umap, jacomy2014forceatlas2, tang2016visualizing, amid2019trimap}. {Our preliminary experiments (see Appendices \ref{app:umap1} and \ref{app:umap2}) suggest that the failure modes discussed in this paper may extend to other force-based dimension reduction techniques.} Our hope is that this work inspires the reader to explore this fascinating landscape further and pursue the essential question: what can be provably deduced from a visualization?



%Our hope is that this work inspires the reader to explore this fertile frontier further and quest after the quintessential question: what can be provably deduced from a visualization?

%for future work is to systematically chart this landscape. to explore this landscape further.

%whose false positives remain largely unstudied, both in theory and practice.

%we would like to provide a taxonomy to understand this zoo of modern techniques.

% Our broader vision is to taxonomize the array 
 
 %understand the intrinsic strengths and limitations of 

%s compared to other methods;–––the gold standard here is separation, not just non sequitir performance guarantee
 
 %taxonomize of other data visualization techniques. Can 

% dissect the favorable and unfavorable qualities


%t-SNE is one of many visualization methods that, in practice, rely purely on a gradient-based optimization of the output points. Related methods in this regard include UMAP \citep{mcinnes2018umap}, ForceAtlas \citep{jacomy2014forceatlas2}, LargeVis \citep{tang2016visualizing}, and TriMap \citep{amid2019trimap}. This is in contrast to an earlier generations of data visualization methods, from PCA and classical MDS  to \citet{tenenbaum2000global} and Laplacian Eigenmaps \citet{belkin2003laplacian}, which usually reduce to eigenvalue problems). Future work should focus on understanding this landscape of ``force-based'' visualization methods and their seemingly unreasonable effectiveness.

 
%It is understood on some heuristic level, for instance, that the affinity matrix asymmetry helps alleviate the ``crowding problem.'' Likewise, 

%Likewise, there is reason to believe that additive invariance has 

%The grand vision is an axiomatic approach to data visualization, analogous to impossibility theorem for data visualization. Can we rigorously define certain desiderate of a data visualization method a la Arora


%Can the benefits of additive invariance be rigorously understood from a denoising perspective?   %we leave it open to future work to quantify the strengths. 

%FORMALIZE POISON POINTS

%It would be instructive to extend our false positive and outlier type studies to other data visualization methods. 

%\textcolor{red}{why the glitch?} 
%One of the distinctive features of t-SNE that arose in our analysis is the property of additive invariance. While we have explored how this property relates to the exaggeration of cluster structure and the sensitivity of t-SNE to poison points, there is reason to believe that additive invariance is beneficial in the processing of high-dimensional, simplex-like data. Additive invariance data processing removes this high-dimensional signature and zeros in on the interrelations between interpoint distances. %It is worth noting that UMAP \citep{mcinnes2018umap} have a feature similar (but not quite identical) to t-SNE's additive invariance, involving the subtraction of the minimum interpoint distance.    %This is because high-dimensional data tends to be close to a simplex. This is a well-known and easily observed principle in practice, see Figure \ref{fig:NLP_simplex}, as well as in theory, from the perspective of concentration of measure. It is possible that additive-invariance is key to designing good data visualization algorithms. Indeed, it is strikingly prevalent among popular visualization algorithms such as t-SNE, UMAP, LargeVis, Laplacian Eigenmaps, and Diffusion Maps. 
%\todo[inline]{Add discussion of: $\E_{x'\sim N(x,I_d), y'\sim N(y,I_d)}[\lVert x' - y' \rVert_2^2] = \lVert x - y \rVert_2^2 + C $ phenomenon.}
%
%\todo[inline]{Include in one form or another in discussion:
%\begin{theorem}\label{thm:additiveInvariance}
%    Laplacian Eigenmaps, LargeVis, SNE, t-SNE, and diffusion maps with Gaussian kernel all are invariant under additive scaling.
%\end{theorem}}

%Indeed, this property may be useful for other endeavors involving high-dimensional data such as measuring cluster significance. Commonly used measures of cluster significance, such as Silhouette Score, rely on comparing the gap between inter-clusters and intra-cluster distances ``multiplicatively'' (i.e. comparing these distances as a ratio). Therefore, high-dimensional datasets which tend to have close to 1 aspect ratios will not present as highly clustered to these measures even though the data may have a far more statistically significant clustering than similar datasets in low dimensions. Our hope is that comparing inter- and intra-cluster distances ``additively'' would overcome this issue and provide a computationally-friendly way to measure cluster significance that is more aligned with the nature of high-dimensional data.

%[additive invariance paragraph?]



%t-SNE is one of many visualization methods that, in practice, rely purely on a gradient-based optimization of the output points. Related methods in this regard include UMAP \citep{mcinnes2018umap}, ForceAtlas \citep{jacomy2014forceatlas2}, LargeVis \citep{tang2016visualizing}, and TriMap \citep{amid2019trimap}. This is in contrast to an earlier generations of data visualization methods, from PCA and classical MDS  to \citet{tenenbaum2000global} and Laplacian Eigenmaps \citet{belkin2003laplacian}, which usually reduce to eigenvalue problems). Future work should focus on understanding this landscape of ``force-based'' visualization methods and their seemingly unreasonable effectiveness.


%\textcolor{red}{i want a sentence name-dropping all the ``force-based'' methods- trimap \citep{amid2019trimap}- umap \citep{mcinnes2018umap}- forceatlas \citep{jacomy2014forceatlas2}- largevis \citep{tang2016visualizing}}


%%%%%%%%%%%% More generally, what information can be deduced about the input given a visualization? 